\section{Status of the uniqueness question and the remaining programme}\label{sec:status}
\subsection{What is and is not established}
Uniqueness of the cycle is proved in two parameter regions: the large-gain region $k\ge1000\max(1,\mu^3)$ of Theorem~\ref{thm:largegain} and its corollaries, and a right-hand neighbourhood of each Hopf point whose width is not explicit (Lemma~\ref{lem:collar}). Theorem~\ref{thm:wedge} adds a statement that is local in phase space: for small rate ratio in the wedge $\rho-1=O(\mu)$, exactly one cycle lies in a fixed energy ball, and it is attracting. That theorem does not exclude further cycles outside the ball, its threshold $\mu_0$ is not explicit, and its local slow-graph input is established in Lemma~\ref{wed:lem:slowgraph} from a checked persistence theorem.

In the complementary region
\begin{equation}\label{eq:openregion}
\cP_{\rm open}=\bigl\{(\mu,k):\ \kH<k<1000\max(1,\mu^3)\bigr\}\ \setminus\ \{\text{Hopf collars}\}
\end{equation}
the cycle count is not determined by this paper. No statement here excludes a neutral cycle in $\cP_{\rm open}$, and none produces one. The floating-point computations of Section~\ref{sec:evidence} found a single cycle, with $\lambda_c<0$, at every parameter examined; they cover a finite set of parameters and cannot detect a pair of cycles closer than the scan spacing.

\subsection{Equivalent forms of the open statement}
By Corollary~\ref{cor:criterion}, uniqueness at $(\mu,k)\in\cP$ follows if every cycle there has $\lambda_c<0$, and non-uniqueness forces a cycle with $\lambda_c\ge0$. For a given cycle the following are equivalent to $\lambda_c<0$:
\begin{enumerate}[label=(\roman*)]
\item $\avg\eta<\mu$ (Proposition~\ref{prop:multipliers});
\item $\displaystyle\int_{U_{\min}}^{U_{\max}}\bigl(\kappa_{\rm up}(U)-\kappa_{\rm down}(U)\bigr)\,\mathrm dU>0$, where $\kappa=K/\Sg$ is evaluated on the rising and falling arcs (Theorem~\ref{thm:gauges}, Proposition~\ref{prop:kappa});
\item $\avg{\mathfrak r}<A$ for the strong-stable ratio $\mathfrak r$ (Lemma~\ref{lem:strongstable});
\item a fixed sign of the window integral in the rank-one identity (Proposition~\ref{prop:rankone}).
\end{enumerate}
Each form isolates one unknown sign. In (ii) the available convexity inequalities bound $\kappa$ from below on both arcs, whereas a proof needs an upper bound on the falling arc. Near the Hopf point the difference in (ii) is of second order in the amplitude and its integrand changes sign at first order, so any argument must be exact to second order at the equilibrium. Section~\ref{sec:negative} records the certificates that were tested and why each fails. It also proves that a single sign change of the one-turn integral of $\eta-\mu$ along the crest half-axis of the attractor would give uniqueness (Proposition~\ref{neg:prop:signchange}, which rests on Proposition~\ref{prop:crestaxis}). Such a sign change follows if the one-turn mean of $\eta-\mu$ is decreasing in the crest amplitude. That statement held at every parameter tested in floating point and is unproved. The stronger statements that the one-turn integral itself is decreasing, or that the return map on the half-axis is sublinear, fail in floating point at extreme rate ratios.

\subsection{A programme for the remaining region}
The proof of Jones' conjecture for Wright's equation combined an analytic large-parameter regime, an explicit neighbourhood of the Hopf point, and computer-assisted enclosure of all periodic solutions on the compact parameter range in between \cite{Xie93,JLM17,vdBJ18,Jaquette19}. The same architecture applies here, with one additional difficulty: the parameter set is two-dimensional and not compact as $\min(\mu,1/\mu)\to0$. By the reciprocal-rate correspondence it suffices to treat $0<\mu\le1$. The pieces are as follows; none of them is carried out in this paper, and the cost figures are design estimates from floating-point probes.

\begin{enumerate}[label=(\alph*)]
\item \emph{An explicit Hopf collar.} In the variable $r=(\rho-1)(1+\mu)^2/\mu$ the weak multiplier of the bifurcating cycle satisfies $\log m_t\approx-2\pi\varepsilon^3r/(1+\varepsilon^2)$, $\varepsilon=\sqrt\mu/(1+\mu)$, as $r\to0$ (twice the growth rate of Appendix~\ref{app:hopf} over one period), and the collar is expected to have uniform width in $r$. A validated treatment of the bifurcation point in the manner of \cite{vdBJ18} would replace the existential Lemma~\ref{lem:collar}.
\item \emph{A lower and uniform large-gain threshold.} Theorem~\ref{thm:largegain} has $\rho_0(\mu)=1000/(\mu(1+\mu))$. Its binding steps are the first foothold on the up-crossing slope and the estimate of the descending phase, where the only available separation from the equilibrium is $\LamE>0$. A threshold bounded in $\rho$ on $\mu\in[\mu_0,1]$ is needed before a computation can close the gap.
\item \emph{A compact core.} For $\mu\in[\mu_0,1]$ between the collar and the large-gain threshold, the statement to certify on each parameter box is that every cycle is hyperbolic attracting. Every cycle has exactly one crest on the section $\{p=0\}$, all crests lie in an explicit region of that section, the crest map is entrywise monotone there (Proposition~\ref{prop:crestmonotone}) with $\det M<e^{-4\pi}$, and a neighbourhood of $E$ can be excluded after a blow-up. A branch-and-bound over the section, with exclusion of boxes that the return map moves off themselves and a Collatz--Wielandt contraction test on the rest, is estimated at $10^5$ to $10^8$ rigorous return-map evaluations for $\mu\in[0.05,1]$, $\rho\le20$. The estimate depends on overestimation factors that were not measured.
\item \emph{The small-rate-ratio end.} As $\mu\to0$ there are two singular regimes: the near-Hamiltonian wedge $\rho-1=O(\mu)$ of Section~\ref{sec:wedge}, selected at third order, and a relaxation regime at fixed $\rho>1$ in which $m_t$ tends to a limit in $(0,1)$ and $T\sim S(\rho)/\mu$. For $r=(\rho-1)(1+\mu)^2/\mu$ in a compact interval, global uniqueness in the wedge would follow from Theorem~\ref{thm:wedge} together with the exclusion of non-attracting cycles outside the energy ball. Section~\ref{sec:wedge} reduces that exclusion to two a-priori estimates for arbitrary cycles (period at most $C/\mu$ and $p\ge-C$) and a spike inequality that is only sketched; none of the three is proved, and a version of Theorem~\ref{thm:wedge} that is uniform as $r\to0$ is also missing. For the relaxation regime only floating-point evidence and a formal one-dimensional limit map are available. No computation at $\mu\ge\mu_0$ removes the need for an analytic theorem on $0<\mu<\mu_0$.
\end{enumerate}
Without item (d) the programme would give uniqueness for $\mu_0\le\mu\le1/\mu_0$ and all $\rho>1$.

\subsection{Conclusion}
The supporting-plane comparison extends from cycles to all bounded orbits and makes the positive-interior attractor a graph of a convex function (with a disk domain only above Hopf); on it the problem is a planar one with divergence $\eta-\mu$. Every cycle is convex and star-shaped in projection and its weak exponent has the exact forms \eqref{eq:mainformula}. For large gain every cycle is attracting, with explicit constants, and the cycle is unique. Between the Hopf collar and that threshold the conjecture of Smith remains open. The package contains neither an all-parameter uniqueness proof nor a multiple-cycle example. The new arguments were each attacked by an internal adversarial reading and by separately implemented numerical tests, recorded in Section~\ref{sec:evidence}; they have not received review by specialists in the field, and historical priority is not established.
